% GLOB-GEO0004 v1.0.0. Geometry and projections follow geometry_library.json
% and geometry_figure_design_standard.json. Six independent parameters;
% notch_tip_radius = 0.5 * notch_width is derived.
% Build from this directory with SOURCE_DATE_EPOCH=0 and
% pdflatex -halt-on-error -no-shell-escape notched_prism.tex. Convert the PDF
% with pdftocairo -svg notched_prism.pdf notched_prism.svg.
%
% Local frame (geometry_library.json): the reference face ABCD is the plane
% y = 0, and O is the midpoint of its longitudinal edge AD, i.e. O lies at
% longitudinal midspan (x = 0) on the edge where the face y = 0 meets the
% face z = 0. x is parallel to AD, y is normal to ABCD, and z extends from AD
% toward BC. The solid occupies -length/2 <= x <= length/2,
% 0 <= y <= thickness, 0 <= z <= depth. The notch enters from the face
% z = depth and penetrates toward decreasing z by the positive scalar
% notch_depth, so its deepest point lies at z = depth - notch_depth. The
% notch runs through the full thickness. Its centerline lies at
% x = notch_offset (signed), and notch_width is measured along x. The notch
% has parallel sides and ends in a semicircle of radius notch_width/2, so the
% straight sides are notch_depth - notch_tip_radius long, which the library
% requires to be strictly positive. The design standard projects +z upward,
% so the notched face z = depth is at the top of every view.
\documentclass[border=1pt]{standalone}
\usepackage[T1]{fontenc}
\usepackage{lmodern}
\usepackage{tikz}
\usetikzlibrary{calc}
% Shared formatting for four-view CIVITA geometry figures.
\usetikzlibrary{arrows.meta}

\definecolor{CivitaGeometry}{HTML}{242424}
\definecolor{CivitaHidden}{HTML}{555555}
\definecolor{CivitaCenterline}{HTML}{525252}
\definecolor{CivitaConstruction}{HTML}{737373}
\definecolor{CivitaAxes}{HTML}{245B83}
\definecolor{CivitaDimension}{HTML}{202020}
\definecolor{CivitaSurface}{HTML}{ECEFF1}
\definecolor{CivitaText}{HTML}{1B1B1B}
\colorlet{CivitaSideSurface}{CivitaSurface!92!CivitaGeometry}
\colorlet{CivitaEndSurface}{CivitaSurface!48!white}

\newcommand{\CivitaPanelWidth}{180}
\newcommand{\CivitaPanelHeight}{134}
\newcommand{\CivitaColumnGap}{12}
\newcommand{\CivitaRowGap}{14}
\newcommand{\CivitaPanelPadding}{8}

\tikzset{
  civita/figure/.style={x=1pt,y=1pt,line cap=round,line join=round},
  civita/edge/.style={draw=CivitaGeometry,line width=.7pt},
  civita/hidden/.style={draw=CivitaHidden,line width=.45pt,dash pattern=on 2pt off 2pt},
  civita/centerline/.style={draw=CivitaCenterline,line width=.45pt,
    dash pattern=on 4pt off 1.5pt on .6pt off 1.5pt},
  civita/construction/.style={draw=CivitaConstruction,line width=.35pt,dotted},
  CivitaArrow/.tip={Latex[length=4pt,width=3pt]},
  civita/dimension/.style={draw=CivitaDimension,line width=.45pt,
    arrows=CivitaArrow-CivitaArrow},
  civita/extension/.style={draw=CivitaDimension,line width=.35pt},
  civita/dimension extension/.style={civita/extension,shorten <=2pt,shorten >=-2pt},
  civita/axis/.style={draw=CivitaAxes,line width=.65pt,arrows=-CivitaArrow},
  civita/label/.style={text=CivitaText,font=\sffamily\fontsize{8}{9}\selectfont,inner sep=1pt},
  civita/axis label/.style={civita/label,text=CivitaAxes},
  civita/title/.style={text=CivitaText,font=\sffamily\bfseries\fontsize{9}{10}\selectfont,
    inner sep=0pt,anchor=north west},
  civita/front/.style={fill=CivitaSurface},
  civita/side/.style={fill=CivitaSideSurface},
  civita/end/.style={fill=CivitaEndSurface}
}

\newcommand{\CivitaCanvas}{%
  \pgfmathsetmacro{\CivitaFigureWidth}{2*\CivitaPanelWidth+\CivitaColumnGap}%
  \pgfmathsetmacro{\CivitaFigureHeight}{2*\CivitaPanelHeight+\CivitaRowGap}%
  \fill[white] (0,0) rectangle (\CivitaFigureWidth,\CivitaFigureHeight);%
}

% Panel indices follow the design standard, with row one at the top.
\newenvironment{CivitaPanel}[3]{%
  \pgfmathsetmacro{\CivitaPanelX}{(#2-1)*(\CivitaPanelWidth+\CivitaColumnGap)}%
  \pgfmathsetmacro{\CivitaPanelY}{(2-#1)*(\CivitaPanelHeight+\CivitaRowGap)}%
  \begin{scope}[shift={(\CivitaPanelX,\CivitaPanelY)}]%
    \node[civita/title] at (\CivitaPanelPadding,{\CivitaPanelHeight-\CivitaPanelPadding}) {#3};%
}{\end{scope}}

\newcommand{\CivitaOrigin}[3][]{%
  \fill[CivitaAxes] (#2) circle[radius=1.3pt];%
  \node[civita/axis label,anchor=#3,#1] at (#2) {$O$};%
}

\newcommand{\CivitaAxis}[4]{%
  \draw[civita/axis] (#1) -- (#2) node[civita/axis label,anchor=#3] {$#4$};%
}

% Named feature points, dimension level in panel coordinates, and symbolic label.
\newcommand{\CivitaHorizontalDimension}[4]{%
  \coordinate (civita-dimension-level) at (0,#3);%
  \draw[civita/dimension extension] (#1) -- (#1 |- civita-dimension-level);%
  \draw[civita/dimension extension] (#2) -- (#2 |- civita-dimension-level);%
  \draw[civita/dimension] (#1 |- civita-dimension-level) -- (#2 |- civita-dimension-level)
    node[civita/label,midway,below=1.5pt] {#4};%
}

% Named feature points, dimension column, label placement options, and label.
\newcommand{\CivitaVerticalDimension}[5]{%
  \coordinate (civita-dimension-column) at (#3,0);%
  \draw[civita/dimension extension] (#1) -- (civita-dimension-column |- #1);%
  \draw[civita/dimension extension] (#2) -- (civita-dimension-column |- #2);%
  \draw[civita/dimension] (civita-dimension-column |- #1) -- (civita-dimension-column |- #2)
    node[civita/label,midway,#4] {#5};%
}

\newcommand{\CivitaParameter}[1]{\textsf{\detokenize{#1}}}

\pdfinfoomitdate=1

% Illustrative construction dimensions in mm. All printed labels are symbolic.
\newcommand{\NPLength}{180}      % length along x
\newcommand{\NPThickness}{60}    % thickness along y
\newcommand{\NPDepth}{80}        % depth along z
\newcommand{\NPOffset}{40}       % notch_offset, signed x of the notch centerline
\newcommand{\NPNotchDepth}{28}   % notch_depth, positive, from the face z = depth
\newcommand{\NPNotchWidth}{12}   % notch_width, opening width along x

% Derived parameter of the library: notch_tip_radius = 0.5 * notch_width.
\pgfmathsetmacro{\NPTipRadius}{0.5*\NPNotchWidth}
% Derived construction values (not parameters of the library).
\pgfmathsetmacro{\NPHalfLength}{\NPLength/2}
\pgfmathsetmacro{\NPxl}{\NPOffset-\NPNotchWidth/2}   % notch wall nearer -x
\pgfmathsetmacro{\NPxr}{\NPOffset+\NPNotchWidth/2}   % notch wall nearer +x
\pgfmathsetmacro{\NPzb}{\NPDepth-\NPNotchDepth}      % notch tip level
\pgfmathsetmacro{\NPzc}{\NPzb+\NPTipRadius}          % tip center, end of the straight sides

% Constraints of geometry_library.json (strict positivity of length,
% thickness, depth, notch_depth, notch_width; notch_depth > 0.5*notch_width,
% so the straight-sided part of the notch is strictly positive;
% notch_depth < depth; the notch lies within the length on both sides).
\newcommand{\NPRequire}[2]{%
  \pgfmathparse{(#1)?1:0}%
  \ifnum\pgfmathresult=0
    \PackageError{civita-notched-prism}{Violated constraint: #2}{Check the parameter values.}%
  \fi}
\NPRequire{\NPLength>0 && \NPThickness>0 && \NPDepth>0}{length, thickness, depth > 0}
\NPRequire{\NPNotchDepth>0 && \NPNotchWidth>0}{notch_depth, notch_width > 0}
\NPRequire{\NPNotchDepth>0.5*\NPNotchWidth}{notch_depth > 0.5 * notch_width}
\NPRequire{\NPNotchDepth<\NPDepth}{notch_depth < depth}
\NPRequire{\NPOffset+0.5*\NPNotchWidth<0.5*\NPLength}{notch_offset + 0.5 * notch_width < 0.5 * length}
\NPRequire{0.5*\NPNotchWidth-\NPOffset<0.5*\NPLength}{0.5 * notch_width - notch_offset < 0.5 * length}

% Notch profile in the x-z plane at a given y: parallel sides at
% x = notch_offset -/+ notch_width/2 from the notched face down to the tip
% center z = zc, then the lower semicircle of radius notch_tip_radius about
% (notch_offset, zc). Its deepest point lies at z = depth - notch_depth.
% \NPArc{y}{a}{b} is the cubic arc from polar angle a to b (degrees, in the
% x-z plane about the tip center) with control length (4/3) tan((b-a)/4) r:
% the standard cubic circle representation, preserved by the affine
% projections used in every panel.
\newcommand{\NPArc}[3]{%
  ({\NPOffset+\NPTipRadius*cos(#2)},#1,{\NPzc+\NPTipRadius*sin(#2)})
  .. controls ({\NPOffset+\NPTipRadius*(cos(#2)-4/3*tan((#3-#2)/4)*sin(#2))},#1,
      {\NPzc+\NPTipRadius*(sin(#2)+4/3*tan((#3-#2)/4)*cos(#2))})
  and ({\NPOffset+\NPTipRadius*(cos(#3)+4/3*tan((#3-#2)/4)*sin(#3))},#1,
      {\NPzc+\NPTipRadius*(sin(#3)-4/3*tan((#3-#2)/4)*cos(#3))})
  .. ({\NPOffset+\NPTipRadius*cos(#3)},#1,{\NPzc+\NPTipRadius*sin(#3)})}
\newcommand{\NPTip}[1]{%
  (\NPxl,#1,\NPzc) -- \NPArc{#1}{180}{270} -- \NPArc{#1}{270}{360}}
\newcommand{\NPTipReversed}[1]{%
  (\NPxr,#1,\NPzc) -- \NPArc{#1}{360}{270} -- \NPArc{#1}{270}{180}}
% Notched x-z outline of the specimen at a given y.
\newcommand{\NPOutline}[1]{%
  (-\NPHalfLength,#1,0) -- (\NPHalfLength,#1,0) -- (\NPHalfLength,#1,\NPDepth)
  -- (\NPxr,#1,\NPDepth) -- \NPTipReversed{#1} -- (\NPxl,#1,\NPDepth)
  -- (-\NPHalfLength,#1,\NPDepth) -- cycle}
% Notch opening region at a given y.
\newcommand{\NPOpening}[1]{(\NPxl,#1,\NPDepth) -- \NPTip{#1} -- (\NPxr,#1,\NPDepth) -- cycle}

\colorlet{NPNotchSurface}{CivitaSurface!78!CivitaGeometry}

\newcommand{\NPOrthoScale}{.60}
\newcommand{\NPIsoScale}{.62}
\pgfmathsetmacro{\NPIsoHorizontal}{\NPIsoScale/sqrt(2)}
\pgfmathsetmacro{\NPIsoDown}{\NPIsoScale/sqrt(6)}
\pgfmathsetmacro{\NPIsoUp}{2*\NPIsoDown}
\newcommand{\NPAxisLength}{16}
\newcommand{\NPIsoAxisLength}{28}

\begin{document}
\begin{tikzpicture}[civita/figure]
  \CivitaCanvas

  \begin{CivitaPanel}{1}{1}{(a) 3D View}
    % View along (-1,-1,-1), with positive z projected upward. The visible
    % faces are x = length/2, y = thickness, and the notched face z = depth;
    % they are never occluded and are painted last. Inside the slot the only
    % surfaces facing the viewer are the side x = notch_offset - notch_width/2
    % and the tip arc between polar angles 180 and 315 degrees (where its
    % inward normal has a positive component along (1,1) in x-z). These two
    % never overlap on screen, both lie in front of the back face y = 0, and
    % wherever neither covers an opening the background shows through the
    % slot. No restriction on thickness or notch_depth is needed.
    \begin{scope}[shift={(77,48.5)},x={(-\NPIsoHorizontal pt,-\NPIsoDown pt)},
      y={(\NPIsoHorizontal pt,-\NPIsoDown pt)},z={(0pt,\NPIsoUp pt)}]
      \newcommand{\NPTopOpening}{(\NPxl,0,\NPDepth) -- (\NPxr,0,\NPDepth)
        -- (\NPxr,\NPThickness,\NPDepth) -- (\NPxl,\NPThickness,\NPDepth) -- cycle}
      \newcommand{\NPBackProfile}{(\NPxl,0,\NPDepth) -- \NPTip{0} -- (\NPxr,0,\NPDepth)}
      \path[fill=NPNotchSurface] (\NPxl,0,\NPzc) -- (\NPxl,\NPThickness,\NPzc)
        -- (\NPxl,\NPThickness,\NPDepth) -- (\NPxl,0,\NPDepth) -- cycle
        % one fill for both surfaces, so no seam appears at their tangency
        \NPArc{\NPThickness}{180}{270} -- \NPArc{\NPThickness}{270}{315}
        -- \NPArc{0}{315}{270} -- \NPArc{0}{270}{180} -- cycle;
      % The notch outline on the back face is visible only through the openings.
      \begin{scope}
        \clip \NPTopOpening;
        \draw[civita/edge] \NPBackProfile;
      \end{scope}
      \begin{scope}
        \clip \NPOpening{\NPThickness};
        \draw[civita/edge] \NPBackProfile;
      \end{scope}
      \path[civita/front,civita/edge] \NPOutline{\NPThickness};
      \path[civita/side,civita/edge] (\NPHalfLength,0,0) -- (\NPHalfLength,\NPThickness,0)
        -- (\NPHalfLength,\NPThickness,\NPDepth) -- (\NPHalfLength,0,\NPDepth) -- cycle;
      \path[civita/end,civita/edge] (-\NPHalfLength,0,\NPDepth) -- (\NPxl,0,\NPDepth)
        -- (\NPxl,\NPThickness,\NPDepth) -- (-\NPHalfLength,\NPThickness,\NPDepth) -- cycle;
      \path[civita/end,civita/edge] (\NPxr,0,\NPDepth) -- (\NPHalfLength,0,\NPDepth)
        -- (\NPHalfLength,\NPThickness,\NPDepth) -- (\NPxr,\NPThickness,\NPDepth) -- cycle;
      % Hidden edges: the rest of the back-face notch outline and the three
      % edges at the occluded corner (-length/2, 0, 0).
      % (overlay: the auxiliary clip rectangle must not enlarge the bounding box)
      \begin{scope}[even odd rule,overlay]
        \clip ([shift={(-200pt,-200pt)}]0,0,0) rectangle ([shift={(200pt,200pt)}]0,0,0)
          \NPTopOpening \NPOpening{\NPThickness};
        \draw[civita/hidden] \NPBackProfile;
      \end{scope}
      \draw[civita/hidden] (-\NPHalfLength,0,0) -- (\NPHalfLength,0,0)
        (-\NPHalfLength,0,0) -- (-\NPHalfLength,\NPThickness,0)
        (-\NPHalfLength,0,0) -- (-\NPHalfLength,0,\NPDepth);
      \coordinate (iso-o) at (0,0,0);
      \coordinate (iso-x) at (\NPIsoAxisLength,0,0);
      \coordinate (iso-y) at (0,\NPIsoAxisLength,0);
      \coordinate (iso-z) at (0,0,\NPIsoAxisLength);
    \end{scope}
    \CivitaAxis{iso-o}{iso-x}{north west}{x}
    \CivitaAxis{iso-o}{iso-y}{north west}{y}
    \CivitaAxis{iso-o}{iso-z}{east}{z}
    \CivitaOrigin[yshift=-4pt]{iso-o}{north}
  \end{CivitaPanel}

  \begin{CivitaPanel}{1}{2}{(b) x-y View}
    % View along -z onto the notched face. The opening spans the thickness.
    \begin{scope}[shift={(84,52)},x={(\NPOrthoScale pt,0pt)},
      y={(0pt,\NPOrthoScale pt)},z={(0pt,0pt)}]
      \draw[civita/edge] (-\NPHalfLength,0) rectangle (\NPHalfLength,\NPThickness);
      \draw[civita/edge] (\NPxl,0) -- (\NPxl,\NPThickness) (\NPxr,0) -- (\NPxr,\NPThickness);
      \coordinate (xy-o) at (0,0);
      \coordinate (xy-x) at (\NPAxisLength,0);
      \coordinate (xy-y) at (0,\NPAxisLength);
      \coordinate (xy-left) at (-\NPHalfLength,0);
      \coordinate (xy-right) at (\NPHalfLength,0);
      \coordinate (xy-nl) at (\NPxl,\NPThickness);
      \coordinate (xy-nr) at (\NPxr,\NPThickness);
      \coordinate (xy-nc) at (\NPOffset,\NPThickness);
    \end{scope}
    \CivitaHorizontalDimension{xy-left}{xy-right}{40}{\CivitaParameter{length}}
    % notch_width: small opening, arrows placed outside the extension lines.
    \coordinate (xy-w-level) at (0,98);
    \draw[civita/dimension extension] (xy-nl) -- (xy-nl |- xy-w-level);
    \draw[civita/dimension extension] (xy-nr) -- (xy-nr |- xy-w-level);
    \draw[civita/extension] ([xshift=-8pt]xy-nl |- xy-w-level) -- ([xshift=8pt]xy-nr |- xy-w-level);
    \draw[civita/dimension,arrows=-CivitaArrow] ([xshift=-8pt]xy-nl |- xy-w-level) -- (xy-nl |- xy-w-level);
    \draw[civita/dimension,arrows=-CivitaArrow] ([xshift=8pt]xy-nr |- xy-w-level) -- (xy-nr |- xy-w-level);
    \node[civita/label,above=1.5pt] at (xy-nc |- xy-w-level) {\CivitaParameter{notch_width}};
    \CivitaAxis{xy-o}{xy-x}{south west}{x}
    \CivitaAxis{xy-o}{xy-y}{west}{y}
    \CivitaOrigin{xy-o}{north east}
  \end{CivitaPanel}

  \begin{CivitaPanel}{2}{1}{(c) x-z View}
    % View along +y. The through-thickness notch is seen in true profile.
    \begin{scope}[shift={(84,32)},x={(\NPOrthoScale pt,0pt)},
      y={(0pt,0pt)},z={(0pt,\NPOrthoScale pt)}]
      \draw[civita/edge] \NPOutline{0};
      \coordinate (xz-o) at (0,0,0);
      \coordinate (xz-x) at (\NPAxisLength,0,0);
      \coordinate (xz-z) at (0,0,\NPAxisLength);
      \coordinate (xz-left) at (-\NPHalfLength,0,0);
      \coordinate (xz-left-top) at (-\NPHalfLength,0,\NPDepth);
      \coordinate (xz-tip) at (\NPOffset,0,\NPzb);
      \coordinate (xz-tip-center) at (\NPOffset,0,\NPzc);
      \coordinate (xz-notch-face) at (\NPOffset,0,\NPDepth);
      \coordinate (xz-depth-col) at ({\NPxl-8},0,0);
    \end{scope}
    % Midspan and notch centerlines, run down to the notch_offset dimension.
    \coordinate (xz-offset-level) at (0,20);
    \draw[civita/centerline] (xz-o |- xz-offset-level) -- ([yshift=6pt]xz-o |- xz-notch-face);
    \draw[civita/centerline] (xz-tip |- xz-offset-level) -- (xz-tip);
    % notch_offset is a signed x-coordinate: measured from midspan, the
    % arrow points to the notch centerline in the positive x direction.
    \draw[civita/dimension,arrows=-CivitaArrow] (xz-o |- xz-offset-level) -- (xz-tip |- xz-offset-level)
      node[civita/label,midway,below=1.5pt] {\CivitaParameter{notch_offset}};
    \CivitaVerticalDimension{xz-left}{xz-left-top}{19}{left=1pt,rotate=90,anchor=south}{\CivitaParameter{depth}}
    % notch_depth: from the notched face inward to the tip, with a leader.
    \draw[civita/dimension extension] (xz-tip) -- (xz-depth-col |- xz-tip);
    % Spans under 10 pt take their arrowheads outside the extension lines.
    \pgfmathparse{\NPNotchDepth*\NPOrthoScale<10?1:0}
    \ifnum\pgfmathresult=0
      \draw[civita/dimension] (xz-depth-col |- xz-tip) -- (xz-depth-col |- xz-notch-face);
    \else
      \draw[civita/extension] (xz-depth-col |- xz-tip) -- (xz-depth-col |- xz-notch-face);
      \draw[civita/dimension,arrows=-CivitaArrow] ([yshift=-7pt]xz-depth-col |- xz-tip) -- (xz-depth-col |- xz-tip);
      \draw[civita/dimension,arrows=-CivitaArrow] ([yshift=7pt]xz-depth-col |- xz-notch-face) -- (xz-depth-col |- xz-notch-face);
    \fi
    \draw[civita/extension] ($(xz-depth-col |- xz-tip)!.5!(xz-depth-col |- xz-notch-face)$)
      -- (xz-depth-col |- 0,92) -- ([xshift=-3pt]xz-depth-col |- 0,92);
    \node[civita/label,anchor=east] at ([xshift=-3.5pt]xz-depth-col |- 0,92) {\CivitaParameter{notch_depth}};
    % Derived tip radius: radial leader through the tip center to the arc,
    % annotated with the library relation in smaller, secondary type.
    \draw[civita/dimension,arrows=-CivitaArrow] (xz-tip |- 0,104) -- (xz-tip);
    \draw[civita/extension] (xz-tip |- 0,104) -- ([xshift=3pt]xz-tip |- 0,104);
    \node[civita/label,font=\sffamily\fontsize{7}{8}\selectfont,align=left,anchor=west]
      at ([xshift=3.5pt]xz-tip |- 0,104)
      {\CivitaParameter{notch_width} =\\ 2\,\CivitaParameter{notch_tip_radius}};
    \CivitaAxis{xz-o}{xz-x}{south west}{x}
    \CivitaAxis{xz-o}{xz-z}{west}{z}
    \CivitaOrigin[xshift=-1.5pt]{xz-o}{north east}
  \end{CivitaPanel}

  \begin{CivitaPanel}{2}{2}{(d) y-z View}
    % View along -x. The notch is internal here; its tip is a hidden line.
    \begin{scope}[shift={(75,32)},x={(0pt,0pt)},
      y={(\NPOrthoScale pt,0pt)},z={(0pt,\NPOrthoScale pt)}]
      \draw[civita/edge] (0,0,0) rectangle (0,\NPThickness,\NPDepth);
      \draw[civita/hidden] (0,0,\NPzb) -- (0,\NPThickness,\NPzb);
      \coordinate (yz-o) at (0,0,0);
      \coordinate (yz-y) at (0,\NPAxisLength,0);
      \coordinate (yz-z) at (0,0,\NPAxisLength);
      \coordinate (yz-right) at (0,\NPThickness,0);
    \end{scope}
    \CivitaHorizontalDimension{yz-o}{yz-right}{20}{\CivitaParameter{thickness}}
    \CivitaAxis{yz-o}{yz-y}{south west}{y}
    \CivitaAxis{yz-o}{yz-z}{west}{z}
    \CivitaOrigin{yz-o}{north east}
  \end{CivitaPanel}
\end{tikzpicture}
\end{document}
